\documentclass[12pt]{article}
%^ Start of document
\usepackage{times}  %Makes text TimesNewRoman
\usepackage{setspace} %Allows you to set single or double space 
\usepackage{biblatex} %Imports biblatex package
\usepackage{graphicx} % Required for inserting images
\usepackage{authblk}
\usepackage[colorlinks=true, urlcolor=blue, linkcolor=blue]{hyperref}
\usepackage{subfig}
\usepackage{float}
\usepackage{tabularx}
\usepackage{multirow}
\usepackage[paperheight=11in,
   paperwidth=8.5in,
   top=20mm,
   bottom=20mm,
   left=40mm,
   right=40mm]{geometry}
\usepackage{moresize}
\usepackage{tikz}
\usepackage{xcolor}
\usepackage{physics}
\usepackage{amssymb}
\usepackage{mathrsfs}
\usepackage{amsmath}
\usepackage{ragged2e}
\usepackage{comment}
\usepackage{xcolor}
\addbibresource{seminar.bib}
\begin{document}
\singlespacing  %Sets single spacing for whole document
\title{Nanotechnology and Material Science Research - Caliskan}

\author[1]{Zachary Tullier}
\affil[1]{Student\\Physics Department\\ University of Houston - Clear Lake \\Houston, Texas\\U.S}
\date{}
\maketitle


\section{Research Directions and Objectives}
    The presentation focused on advancements in material science and nanotechnology, particularly the use of boron nitride nanotubes (BNNTs) as an alternative to carbon nanotubes (CNTs). The discussion explored using advanced quantum simulation tools to design molecular structures, optimize them, and evaluate their mechanical, thermal, and electronic properties and bridging the gap between theoretical predictions and real-world applications.


\section{Applications Across Disciplines}
    Nanotechnology serves as a multidisciplinary field connecting physics, biology, chemistry, engineering, and information technology. Highlighted applications included:
    \begin{itemize}
        \item Nanowires, Nanotubes, and Molecular Junctions
        \item Diodes and Transistors
        \item Integrating biological systems with spintronic technology.
        \item Creating spin-polarized systems for data storage and processing.
        \item Using nanomaterials for medical implants or treatments.
        \item Employing nanotubes as detectors for biological or chemical changes.
        \item Developing materials capable of efficient hydrogen adsorption for fuel cells.
        \item Utilizing high thermal conductivity for heat dissipation.
    \end{itemize}

\section{Boron Nitride Nanotubes (BNNTs)}
    BNNTs were highlighted in this presentation. BNNTs have a wide band gap (~5.5 eV), classifying them as insulators. However, this band gap can be altered through doping or structural modifications. The band gap plays a crucial role in determining whether the material behaves as a semiconductor or an insulator. BNNTs are mechanically robust, with high Young’s modulus values indicative of significant strength and resistance to deformation. They exhibit superior thermal conductivity, essential for applications requiring efficient heat transfer. BNNTs can be categorized into \textbf{zigzag} and \textbf{armchair} structures, with differences in their electronic and mechanical properties.

    The presenter emphasized BNNTs as an alternative to CNTs due to their consistent semiconducting properties, high strength, and thermal advantages.

\section{Role of Transition Metal Doping}
    The introduction of transition metals, such as nickel and cobalt, into BNNTs enhances their functionality. Pure BNNTs exhibit no magnetic moments. Doping introduces magnetic properties, making them suitable for spintronic applications. Doped BNNTs exhibit spin polarization, a key requirement for spintronic devices. Despite doping, BNNTs retain their mechanical strength, ensuring structural reliability for practical applications.
    

    The study demonstrated that transition metal doping creates a balance between maintaining mechanical properties and enhancing magnetic functionality.

\section{Comparative Analysis: BNNTs vs. Carbon Nanotubes (CNTs)}
    The presenter compared BNNTs with CNTs, discussing their respective advantages and challenges:
    While CNTs sometimes achieve higher maximum values, BNNTs provide more consistent and reliable thermal properties. CNTs can be metallic or semiconducting based on their chirality, making them less predictable. BNNTs are inherently semiconducting, with properties modifiable by doping. BNNTs were favored for their insulating nature and stability, while CNTs were noted for their variability and potential in other niche applications.

    CNTs are currently the most popular nanotube material, but research is highlighting the benifit of other options. 
    
\section{Computational and Theoretical Approaches}
    The research relied on quantum simulations to model and predict material behavior. The Quantum ATK Software is a powerful tool for designing and testing molecular structures, optimizing parameters, and visualizing results.
    
    \textbf{Key Metrics Analyzed}:
    \begin{itemize}
        \item \textbf{Energy Band Gap}: Indicates electronic behavior.
        \item \textbf{Magnetic Moments}: Reflects the material's suitability for spintronic applications.
        \item \textbf{Density of States (DOS)}: Shows the distribution of energy levels, crucial for understanding electronic properties.
        \item \textbf{Formation Energy}: Indicates structural stability; negative values suggest stable configurations.
    \end{itemize}

    Computational approaches provided valuable insights, guiding experimental setups and real-world validations.

\section{Future Applications and Directions}
    The discussion concluded with insights into future applications of BNNTs. BNNTs doped with transition metals can be used as spin filters, spin field-effect transistors, and other spin-based devices. Hydrogen storage capabilities make BNNTs attractive for renewable energy solutions. BNNTs' superior thermal and mechanical properties position them for applications in heat management and structural components.

\section{Conclusion}
This extended discussion showcased the synergy between theoretical and experimental research in advancing material science. Boron nitride nanotubes emerged as a promising material for next-generation technologies due to their robust thermal, mechanical, and electronic properties. By integrating computational tools, the research provided a pathway to develop innovative solutions in spintronics, energy efficiency, and nanotechnology.

%\printbibliography
\end{document}
